\section{Discussion}
\label{sec:discussion}

\textbf{What the evidence supports.} Three claims are supported by ten-seed paired tests on patient-deduplicated data. First, on this benchmark synchronising less does not cost accuracy; every sparse schedule matches or slightly exceeds full-rate personalisation, and \tierfed{} is non-inferior to FedPer on AUROC and AUPRC at a fraction of the upload. Second, timing and message-size reductions compose multiplicatively when compressors act on the update with error feedback, giving more than two orders of magnitude less upload at unchanged accuracy. Third, under client-level DP the release count is a design variable that must be fixed before training: a schedule that reacts to client updates is itself a leak, and once the schedule is planned, the period trades drift against the noise each release must carry. At matched $\varepsilon$, the planned schedule \ifflag{dpmQHalfEpsFourTfdpVsFedperBetter}{retains more utility than full-rate personalisation}{does not significantly outperform full-rate personalisation} at $\varepsilon=4$, $q=0.5$. A communication quota and a privacy budget therefore have to be negotiated together, which matters operationally because they are usually set by different stakeholders.

\textbf{What it does not support.} \tierfed{} is not more accurate than simpler round-sparse alternatives. A fixed uniform period of four rounds is the most accurate configuration we ran, and AdaComm reaches comparable accuracy with fewer bytes when DP is off. The divergence-driven mechanism that motivates per-tier intervals rarely activates on this data; at most quotas the exact solver returns the same interval for both shared tiers. Under DP, neither adaptive \tierfed{} nor AdaComm is eligible at all, and \tierfed{}-DP with a flat prior always chooses equal intervals for both tiers, so it is a privacy-aware period planner rather than a per-tier scheduler. The case for \tierfed{} over these alternatives is operational rather than statistical: it exposes a quota that maps directly onto a slice or window allocation, solves for the schedule exactly, and composes with compression and DP accounting, whereas AdaComm exposes a loss-driven period and a fixed schedule exposes nothing an operator can set from a byte budget.

\textbf{A retraction.} An earlier version of this work reported that fixed multi-rate schedules lose about 0.05~AUROC to personalised baselines. With one encounter per patient, patient-grouped splits and the fallback-tier bug removed, that result does not reproduce. Readers who relied on it should not.

\textbf{Limitations.}
(i) The sites of the main tabular benchmark are specialty buckets with a share cap, not hospitals; the real-institution federations are small (four heart-disease centres; the open eICU demo hospitals), and the full credentialed eICU-CRD was not run.
(ii) The bound is for SGD and the main benchmark uses Adam; the SGD arm shows that the ordering of methods depends on the local solver.
(iii) $\hat\Gamma_g$ is a dispersion proxy that scales with the interval that produced it; the controller is a feedback loop on that proxy whose stability we observe but do not prove.
(iv) The network model and the socket emulation omit radio scheduling, fading, loss and the downlink; the emulation is application-layer, not kernel-level, and the profiles are representative targets rather than measurements.
(v) The inversion attack is generic, and client-level budgets of $\varepsilon\ge1$ over 20 sites are weak guarantees.
(vi) Poisson participation is used only in the DP experiments; the bound (\ref{eq:bound}) is stated for full participation and does not include a sampling term.
(vii) The FT-Transformer arm was trained with a non-deterministic GPU attention kernel; all other results are bit-reproducible.
(viii) On the imaging benchmark under strong label skew, the methods that share the classifier head are the most accurate and both head-private methods trail them; the results here do not transfer to that setting.
(ix) Ten seeds on one public tabular dataset and one imaging collection is a benchmark, not a clinical validation.

\textbf{Practical guidance.} For a tabular model of this size, an operator who needs a byte budget should set $\rho$ from the transfer window through (\ref{eq:compose}) and (\ref{eq:netcost}), compose the schedule with 8-bit quantisation or sparse ternary compression with error feedback, and, when client-level DP is required, fix the schedule before training (a fixed period or \tierfed{}-DP), calibrate $\sigma$ to its release count, and never use a scheduler that reacts to client updates.
